在模 6 的世界里 \(2\times3=0\)，两个非零的东西乘出了零，于是``两边同除''这个从小学起就没出过错的动作忽然不能用了。另一头，编译器碰到两个分支分别返回 \texttt{Int32} 与 \texttt{Float64} 时，得为整个表达式选一个类型：既要容得下两支，又不能宽到让后续检查失效，它要的是所有共同上界里最小的那个。

这两件事看起来不相干，一个属于代数，一个属于次序，但它们问的是同一个问题：手上这套对象，究竟保留了哪些运算规律？答案决定了哪些熟悉的动作还能做。上一章的群只管一种运算，够用来讨论对称，不够用来讨论``加和乘怎么配合''。本章沿两条路线各走一步。

\begin{center}
\begin{tikzpicture}[x=1cm,y=1cm,font=\footnotesize,>=stealth]
  \draw[black!50,fill=black!4] (3.6,4.4) rectangle (8.6,5.2);
  \node[font=\small] at (6.1,4.8) {一个集合，你决定保留哪些规律};

  \draw[->,black!60] (5.0,4.35) -- (2.8,3.55);
  \node[black!70,above,rotate=20] at (3.7,3.8) {叠运算};
  \draw[->,black!60] (7.2,4.35) -- (9.4,3.55);
  \node[black!70,above,rotate=-20] at (8.5,3.8) {叠次序};

  \draw[black!55,fill=orange!10] (1.1,2.7) rectangle (3.6,3.5);
  \node at (2.35,3.1) {环};
  \draw[black!55,fill=orange!16] (1.1,1.4) rectangle (3.6,2.2);
  \node at (2.35,1.8) {整环};
  \draw[black!55,fill=orange!24] (1.1,0.1) rectangle (3.6,0.9);
  \node at (2.35,0.5) {域};
  \draw[->,black!60] (2.35,2.65) -- (2.35,2.25);
  \draw[->,black!60] (2.35,1.35) -- (2.35,0.95);
  \node[right,black!65,align=left] at (3.75,2.45) {无零因子};
  \node[right,black!65,align=left] at (3.75,1.15) {非零元可逆};

  \draw[black!55,fill=blue!10] (8.6,2.7) rectangle (11.1,3.5);
  \node at (9.85,3.1) {偏序};
  \draw[black!55,fill=blue!16] (8.6,1.4) rectangle (11.1,2.2);
  \node at (9.85,1.8) {格};
  \draw[black!55,fill=blue!24] (8.6,0.1) rectangle (11.1,0.9);
  \node at (9.85,0.5) {完全格};
  \draw[->,black!60] (9.85,2.65) -- (9.85,2.25);
  \draw[->,black!60] (9.85,1.35) -- (9.85,0.95);
  \node[left,black!65,align=right] at (8.45,2.45) {两元素有上下确界};
  \node[left,black!65,align=right] at (8.45,1.15) {任意子集有上下确界};
\end{tikzpicture}

\emph{两条路线，同一种兑换方式：往下一格多要求一条性质，换回一样能做的事。左边换来的是消去与除法，右边换来的是``合并''成为运算、以及不动点必定存在。反过来读也成立——掉出某一格，就会丢掉对应的那样能力。}
\end{center}

第一篇沿左路走。环把加法和乘法装进同一个结构，加法那半照搬交换群，乘法那半只要结合，两者靠分配律通信。省下``非零元可逆''这条要求之后，除法失效、零因子登场，模合数运算里那些反直觉的现象、有限域为什么必须用不可约多项式来造，都从这一处解释得通。哈希、编码与密码学的字节运算都发生在有限域上，这一篇是它们的地基。

第二篇沿右路走。偏序只回答``谁在谁下面''，一旦保证任意两元素的上下确界存在，``取上界''与``取下界''就升格成两种可以连写的运算：集合的交并、\(\gcd\) 与 \(\operatorname{lcm}\)、类型系统里的子类型合并、权限模型里的能力叠加，全是它的化身。篇末的 Knaster--Tarski 定理说明完全格上的单调映射必有最小不动点——递归定义与程序分析中的``最小解''之所以存在，依据就在这里。

第三篇把右路的抽象用出去。Galois 连接是两个偏序集之间一对方向相反的映射，一边放松约束、一边收紧结论，往返一次的复合是闭包算子。抽象解释、约束传播、形式概念分析共用这一副骨架；而``迭代到收敛''这类算法为什么一定停得下来，答案落在格的高度上。

\begin{itemize}
\tightlist
\item \hyperref[sec:04-Discrete-Mathematics/08-Rings-and-Lattices/01]{``环把两种运算放进同一结构''}：判断一个代数动作在当前对象上合不合法，先看零因子与可逆元；
\item \hyperref[sec:04-Discrete-Mathematics/08-Rings-and-Lattices/02]{``格把偏序变成代数运算''}：判断一个合并操作是不是 join，别看 Hasse 图画得漂不漂亮；
\item \hyperref[sec:04-Discrete-Mathematics/08-Rings-and-Lattices/03]{``格在语义、约束与 Galois 连接中的角色''}：判断一次近似是否可靠，看两个方向的映射是否配成一对。
\end{itemize}

两条路线彼此独立，读格不需要先学环，读环也用不上偏序，从哪一篇进都行。若从模运算的困惑进来，可以先补\hyperref[ch:029]{``初等数论：整除结构怎样变成可计算算法''}里的模逆；若从逻辑进来，\hyperref[ch:033]{``布尔代数与自动机：有限规则怎样处理逻辑与序列''}里的 AND 与 OR 就是分配格上的 meet 与 join，但别把分配律和补元当成所有格都自带的东西。第一遍读完能回答三个问题就够了：手上这个非零元素可逆吗？这堆上界里有没有最小的那一个？这个迭代凭什么会停？
